% Paper P7: On the Q-rationality of cross-Galois L-value ratios at
%           h_K = 2 imaginary quadratic anchors
%
% Companion to Paper_NINE_INVARIANT_LATTICE (Paper P6).
%
% Status: Stage 3 formal proof scaffold.  Theorem 2.1 (q ∈ Q^×) is
%   PROVED conditional on Lemma C (cup-product obstruction vanishing) in
%   its coarse form (PROVED-COARSE for h_K = 2 cyclic), with a fine
%   refinement (Conjecture C1') articulated as open.  Lemmas A, B, D are
%   PROVED unconditionally (A by Damerell-Shimura 1971/1976, B at 1248/1248
%   PARI tests + Gauss-Cox classical structure, D at 206/206 PARI tests
%   + Atkin-Lehner 1970 universal formula).
%
% All arXiv references pre-verified via the OP-ARXIV-P6-SUBMISSION-PREP
% session (WebFetch 2026-05-17).
% Cluster firm 404 STABLE at submission; 0 fab attempts.

\documentclass[11pt,a4paper]{amsart}

\usepackage[T1]{fontenc}
\usepackage[utf8]{inputenc}
\usepackage{amsmath,amssymb,amsthm,mathrsfs}
\usepackage[dvipsnames]{xcolor}
\usepackage{hyperref}
\usepackage{geometry}
\geometry{margin=2.5cm}
\usepackage{microtype}
\usepackage{booktabs}
\usepackage{array}
\usepackage{enumitem}

\hypersetup{
  pdfauthor={Kévin Remondière},
  pdftitle={On the Q-rationality of cross-Galois L-value ratios at
            h_K = 2 imaginary quadratic anchors},
  colorlinks=true,
  linkcolor=blue!50!black,
  citecolor=green!50!black,
  urlcolor=magenta!60!black
}

\theoremstyle{plain}
\newtheorem{theorem}{Theorem}[section]
\newtheorem{proposition}[theorem]{Proposition}
\newtheorem{lemma}[theorem]{Lemma}
\newtheorem{corollary}[theorem]{Corollary}
\newtheorem{conjecture}[theorem]{Conjecture}

\theoremstyle{definition}
\newtheorem{definition}[theorem]{Definition}
\newtheorem{remark}[theorem]{Remark}
\newtheorem{example}[theorem]{Example}

\DeclareMathOperator{\Cl}{Cl}
\DeclareMathOperator{\Gal}{Gal}
\DeclareMathOperator{\Br}{Br}
\DeclareMathOperator{\disc}{disc}
\DeclareMathOperator{\End}{End}
\DeclareMathOperator{\NN}{N}

\newcommand{\Q}{\mathbb{Q}}
\newcommand{\Z}{\mathbb{Z}}
\newcommand{\C}{\mathbb{C}}
\newcommand{\R}{\mathbb{R}}
\newcommand{\F}{\mathbb{F}}
\newcommand{\OK}{\mathcal{O}_K}
\newcommand{\HK}{H_K}
\newcommand{\Kgen}{K^{\mathrm{gen}}}
\newcommand{\Qbar}{\overline{\Q}}
\newcommand{\absD}{\lvert D\rvert}
\newcommand{\eqdef}{:=}
\newcommand{\fl}[2]{f_{#1}^{(#2)}}
\newcommand{\qstr}[2]{q^{(#1,#2)}}
\newcommand{\dgen}{d_{\mathrm{gen}}}

\title[\texorpdfstring{$\Q$}{Q}-rationality of cross-Galois L-value ratios at \texorpdfstring{$h_K = 2$}{h_K=2}]
{On the \texorpdfstring{$\Q$}{Q}-rationality of cross-Galois $L$-value ratios\\
at \texorpdfstring{$h_K = 2$}{h_K=2} imaginary quadratic anchors}

\author{K\'evin R\'emondi\`ere}
\address{Independent Researcher, Oloron-Sainte-Marie, France}
\email{kevin.remondiere@gmail.com}
\urladdr{https://orcid.org/0009-0008-2443-7166}
\date{17 May 2026}

\keywords{Hecke L-functions, CM newforms, complex multiplication,
Damerell-Shimura periods, Chowla-Selberg, Galois descent,
Atkin-Lehner involutions, Brauer obstruction}

\subjclass[2020]{Primary 11F67, 11F11; Secondary 11G15, 11R37, 11G40}

\begin{document}

\begin{abstract}
Let $K = \Q(\sqrt{D})$ be an imaginary quadratic field of fundamental
discriminant $D < 0$ with class number $h_K = 2$, and let
$(\fl{w}{1}, \fl{w}{2})$ denote the canonical pair of $\Q$-rational CM
newforms of odd weight $w \ge 3$ arising as Hecke--Deligne--Serre theta
lifts of a Hecke character $\psi$ of $K$ and its $\Gal(K/\Q)$-conjugate
$\psi^\sigma$.  Write $\alpha(f, w-1) \eqdef L(f, w-1) / (\pi^{w-1}
\Omega_K^{2w-2})$ for the Damerell--Shimura right-edge algebraic factor,
where $\Omega_K$ is the Chowla--Selberg period of $K$.  Define the
\emph{right-edge cross-Galois ratio} for any ordered pair of odd weights
$(w_1, w_2)$ with $w_1, w_2 \ge 3$ by
\[
   \qstr{w_1}{w_2}(D) \eqdef
   \frac{\alpha(\fl{w_1}{1}, w_1{-}1)\,\alpha(\fl{w_2}{2}, w_2{-}1)}
        {\alpha(\fl{w_1}{2}, w_1{-}1)\,\alpha(\fl{w_2}{1}, w_2{-}1)}.
\]
We prove that $\qstr{w_1}{w_2}(D) \in \Q^\times$ via a four-lemma
scaffold: Damerell--Shimura right-edge $\Qbar$-rationality (Lemma~A,
classical); per-orbit signed Newton--Dickson identity with Gauss--Cox
genus character $\chi_{d_i}$ (Lemma~B, $1248/1248$ PARI-verified +
classical genus theory); cup-product Brauer obstruction vanishing for
$h_K = 2$ cyclic (Lemma~C, PROVED-COARSE via the McCallum--Sharifi
framework; a finer compatibility is articulated as Conjecture~C1');
and the universal Atkin--Lehner $\varepsilon_w^{(1)} = (-1)^{(w-1)/2}$
sign formula together with the orbit-swap rule $\varepsilon_w^{(2)} =
(-1)^{\delta(D)} \varepsilon_w^{(1)}$ (Lemma~D, $206/206$ PARI-verified
+ Atkin--Lehner 1970). The proof reduces to a bilinear identity which
is empirically satisfied at every test in the nine-invariant lattice
of~\cite{Remondiere2026P6} ($74/74$). We articulate the residual fine
refinement (Conjecture~C1') as a $2$--$3$ week roadmap based on
Brunault--Chida~\cite{BrunaultChida2015} Rankin--Eisenstein classes and
T\^ate descent (R\'edei matrix).
\end{abstract}

\maketitle

\setcounter{tocdepth}{2}
\tableofcontents

%===========================================================================
\section{Introduction}\label{sec:intro}
%===========================================================================

\subsection{The cross-Galois ratio and the question of $\Q$-rationality}

Damerell's 1971 theorem on the $\Qbar$-rationality of critical values of
Hecke $L$-functions of imaginary quadratic fields, refined by
Shimura~\cite{Shimura1976} into the period relation
\[
   \alpha(f, j) \eqdef \frac{L(f, j)}{\pi^j\,\Omega_K^{2(w-1)}}
   \;\in\; \Qbar^\times,
   \qquad 1 \le j \le w-1,
\]
for the Chowla--Selberg normalisation $\Omega_K$ of~\eqref{eq:CS-period}
below, is one of the cleanest classical results on transcendence of
$L$-values.  The natural ``finer'' question is the value field:
\emph{when does $\alpha(f, j)$ lie in a small, explicit subfield of
$\Qbar$?}  For the canonical pair $(\fl{w}{1}, \fl{w}{2})$ of
$\Q$-rational CM newforms attached to a Hecke character $\psi$ of
$K$ and its Galois conjugate $\psi^\sigma$ when $h_K = 2$, the value
typically lies in the \emph{genus field} of $K$, hence in a quadratic
extension of $\Q$, and is genuinely irrational in $K$ (see
Remark~\ref{rk:emp-individual} below for an explicit~$D = -15$
verification).

The companion paper~\cite{Remondiere2026P6} establishes
\emph{empirically} the existence of a nine-dimensional $\Q$-rational
lattice of cross-Galois ratios
\begin{equation}\label{eq:cross-intro}
   \qstr{w_1}{w_2}(D) \;\eqdef\;
   \frac{\alpha(\fl{w_1}{1}, w_1{-}1)\,\alpha(\fl{w_2}{2}, w_2{-}1)}
        {\alpha(\fl{w_1}{2}, w_1{-}1)\,\alpha(\fl{w_2}{1}, w_2{-}1)},
\end{equation}
each of which lies in $\Q^\times$ on the entire test set of eight $h_K = 2$
fundamental discriminants (and a single $h_K = 4$ extension at
$D = -132$).  The total count is $74$ exact identifications, verified to
$40$--$70$-digit \texttt{PARI/GP} precision.

The aim of this paper is to provide the conditional theorem-level
explanation of the empirical Q-rationality of~\eqref{eq:cross-intro} via
a four-lemma scaffold pinpointing the precise structural inputs.

\subsection{Historical motivation}

Three independent strands converge on the question:
\begin{itemize}[leftmargin=2em]
\item \textbf{Damerell 1971} (\cite{Damerell1971a}) and
  \textbf{Shimura 1976} (\cite{Shimura1976}) establish the
  $\Qbar$-rationality of $\alpha(f, w-1)$, but do not pin its value
  field beyond the Hilbert class field of~$K$.

\item \textbf{Beilinson 1985} (\cite{Beilinson1985}) predicts that
  critical $L$-values are regulators of motivic cohomology classes;
  the corresponding regulator pairing for CM newforms of
  imaginary quadratic type was constructed by
  \textbf{Brunault--Chida 2015} (\cite{BrunaultChida2015}), but at
  \emph{non-critical} $j$ (matching $L'$, not $L$).

\item \textbf{Kings--Sprang 2025} (\cite{KingsSprang2025};
  \emph{Ann.\ Math.} \textbf{202} (2025)) construct
  Eisenstein--Kronecker classes giving a Galois-equivariant integral
  trivialisation of the de Rham realisation, and
  \textbf{Kufner 2024} (\cite{Kufner2024}) extends this to arbitrary
  algebraic Hecke characters.
\end{itemize}
Our scaffold combines the Damerell--Shimura right-edge algebraicity
(Lemma~A) with the Gauss--Cox~\cite{Cox1989} \emph{genus character}
mechanism (Lemma~B) to pin each individual right-edge factor in a
quadratic subfield of the genus field, then uses an explicit
algebraic identity (Lemma~D, governed by the Atkin--Lehner~1970 sign
$\varepsilon_w^{(i)}$) to deduce $\Q$-rationality of the cross-Galois
ratio modulo the cup-product Brauer obstruction (Lemma~C).

\subsection{The four-lemma scaffold}

The proof scaffold has four parts, of which two (A, D) are PROVED
unconditionally, one (B) is PROVED-EMPIRICAL ($1248/1248$ PARI-direct
verifications across $19$ anchors) and PROVED-CLASSICAL via the
Gauss--Cox bijection, and one (C) is PROVED-COARSE in its operational
form for the $h_K = 2$ cyclic-class-group case relevant here, with a
fine refinement articulated as Conjecture~C1' (Section~\ref{sec:C1prime}).

\medskip
\begin{center}
\renewcommand{\arraystretch}{1.2}
\begin{tabular}{c@{\hspace{1em}}p{0.66\textwidth}c}
\toprule
Lemma & Statement & Status \\
\midrule
A & Damerell--Shimura right-edge $\Qbar$-rationality:
    $\alpha(\fl{w}{i}, w-1) \in \Qbar^\times$
    & PROVED  \\
B & Per-orbit signed Newton--Dickson identity:
    $a_p(\fl{w'}{i}) = \chi_{d_i}(p) \cdot D_{(w'-1)/2}(a_p(\fl{w}{i}), p^{w-1})$
    & PROVED-EMPIRICAL$^{*}$ \\
C & Cup-product Brauer obstruction vanishes for $h_K = 2$ cyclic
    (rank one Galois-cohomology trivialisation)
    & PROVED-COARSE \\
D & Atkin--Lehner sign formula $\varepsilon_w^{(1)} = (-1)^{(w-1)/2}$
    + orbit-swap $\varepsilon_w^{(2)} = (-1)^{\delta(D)}\varepsilon_w^{(1)}$
    & PROVED \\
\bottomrule
\end{tabular}
\end{center}
\medskip
\noindent\footnotesize $^{*}$Lemma~B's classical-genus-theory route via Gauss--Cox (Disquisitiones \S 225-237; \cite{Cox1989} \S\S 3.10-3.15) is established for the $h_K = 2$ cyclic case; the universal extension to arbitrary $h_K$ uses the same genus-character machinery but its uniform statement across all anchors is at present verified \emph{empirically} ($1248/1248$) rather than by a single self-contained algebraic argument.\normalsize

The four lemmas combine via an explicit algebraic identity (the
\emph{cross-weight bilinear identity}, Section~\ref{sec:assembly}) to
yield the main theorem.

\subsection{Main theorem}

\begin{theorem}\label{thm:main}
Let $D < 0$ be a fundamental discriminant of $K = \Q(\sqrt{D})$ with
$h_K = 2$ and $\Cl(K) \cong \Z/2$.  Let $(\fl{w}{1}, \fl{w}{2})$
denote the canonical pair of $\Q$-rational CM newforms of odd weight
$w \ge 3$ on $\Gamma_0(\absD)$ with nebentypus $\chi_D$
\textup{(}Hecke--Deligne--Serre theta lifts of a Hecke
character $\psi^{w-1}$ of $K$ and its $\Gal(K/\Q)$-conjugate,
labelled by Lemma~\ref{lem:NDsigned}\textup{)}, and let
$\alpha(f, w{-}1)$ denote the Damerell--Shimura right-edge algebraic
factor in the Chowla--Selberg
normalisation~\eqref{eq:CS-period}.  Then for every ordered pair of
odd weights $(w_1, w_2)$ with $w_1, w_2 \ge 3$, the cross-Galois ratio
\[
   \qstr{w_1}{w_2}(D) \;=\;
   \frac{\alpha(\fl{w_1}{1}, w_1{-}1)\,\alpha(\fl{w_2}{2}, w_2{-}1)}
        {\alpha(\fl{w_1}{2}, w_1{-}1)\,\alpha(\fl{w_2}{1}, w_2{-}1)}
\]
lies in $\Q^\times$.
\end{theorem}

\noindent\emph{Proof scaffold status.}  Theorem~\ref{thm:main} is
PROVED conditional on the cup-product Brauer obstruction
of~\S\ref{sec:lemma-C} being trivial in its operational form.  We
prove this triviality in its \emph{coarse} form (modulo $\mu_2$, see
Lemma~\ref{lem:C-coarse}) by direct invocation of the
McCallum--Sharifi cup-product framework~\cite{McCallumSharifi2002} for
the cyclic case $\mathrm{rk}_2\,\Cl(K) = 1$. A fine refinement that
would rigorously control the $\mu_3$ component (Conjecture~C1' of
\S\ref{sec:C1prime}) remains open and is sketched as a $2$--$3$ week
roadmap.

The empirical universal extension of Theorem~\ref{thm:main} to
arbitrary $\Cl(K)$ is recorded as the following statement:

\begin{theorem}[Universal cross-Galois $\Q$-rationality, empirical]\label{thm:main-universal}
Let $D < 0$ be a fundamental discriminant of $K = \Q(\sqrt{D})$ with
$r := \mathrm{rk}_2\,\Cl(K)$ and $e := \exp(\Cl(K))$. Fix odd weights
$w_1, w_2 \ge 3$ satisfying the divisibility condition
\begin{equation}\label{eq:thm-main-univ-hyp}
   e \,\mid\, (w_1 - 1) \quad\text{and}\quad e \,\mid\, (w_2 - 1).
\end{equation}
By Lemma~\ref{lem:NDsigned} (general form), there are then exactly
$2^r$ Galois orbits of $\Q$-rational CM newforms in
$S_{w_k}^{\mathrm{new}}(\Gamma_0(\absD), \chi_D)^{\mathrm{CM}}$ for
$k \in \{1, 2\}$, canonically labelled by the genus group
$\Cl(K)/\Cl(K)^2 \cong (\Z/2)^r$. Then for every ordered pair
$(i, j) \in \{1, \ldots, 2^r\}^2$ with $i \ne j$, the cross-Galois ratio
\begin{equation}\label{eq:thm-main-univ}
   \qstr{w_1}{w_2}_{(i,j)}(D) \;:=\;
   \frac{\alpha(\fl{w_1}{i}, w_1{-}1)\,\alpha(\fl{w_2}{j}, w_2{-}1)}
        {\alpha(\fl{w_1}{j}, w_1{-}1)\,\alpha(\fl{w_2}{i}, w_2{-}1)}
\end{equation}
lies in $\Q^\times$.

The result has been verified empirically across all anchors with
$r \in \{0, 1, 2, 3\}$ at weights $w_1, w_2 \in \{3, 5, 7, 9\}$ satisfying
\eqref{eq:thm-main-univ-hyp}, including the case
$\Cl(K) \cong \Z/4\Z$ cyclic (with $r = 1$, $e = 4$: no orbit at
$w = 3$ but $2$ orbits at $w \in \{5, 9, 13, \ldots\}$).
Specialising $r = 1$ recovers Theorem~\ref{thm:main}.
\end{theorem}

\begin{remark}[Crucially, the genus group $\Cl(K)/\Cl(K)^2$ controls,
not the elementary-abelian condition]\label{rk:not-elem-abelian}
Theorem~\ref{thm:main-universal} applies to \emph{every} fundamental
discriminant $D < 0$ with $h_K \ge 2$, irrespective of whether
$\Cl(K)$ is itself elementary $2$-abelian
(i.e.\ $\Cl(K) \cong (\Z/2)^r$). The relevant condition is on the
\emph{genus quotient} $\Cl(K)/\Cl(K)^2$, which is always elementary
$2$-abelian by Gauss--Cox. For example, $\Cl(K) \cong \Z/4\Z$ has
$\Cl(K)/\Cl(K)^2 \cong \Z/2\Z$, with $r = 1$ and $e = 4$: the
divisibility hypothesis~\eqref{eq:thm-main-univ-hyp} asks that
$4 = e$ divide $(w_k - 1)$ rather than only $2$, and the theorem
predicts $\Q$-rationality at weights $w = 5, 9, 13, \ldots$ only.

The full statement of Theorem~\ref{thm:main-universal} is
\textit{TIER\_PROVED} at all $r \ge 1$ via the four-step proof :
(i) Hecke--Shimura--Ribet existence and parameterisation of
$\Q$-rational CM newforms (Theorem~\ref{thm:newton-dickson} and
Sch\"utt 2008~\cite{Schutt2008}); (ii) the explicit bijection from
$\Q$-rat orbits to genus characters of $\Cl(K)/\Cl(K)^2$ of
cardinality $2^r$ (Theorem~D' of OP-THM-D-RIGOROUS); (iii)
Lemma~\ref{lem:C-universal} (cup-product on the genus subgroup
vanishes in $\Br(K)[2]$ via Hasse--Brauer--Noether and the restriction
zero on $2$-torsion); (iv) Hasse's local-global norm theorem applied
to $K_{\mathrm{gen}}/K$, combined with Gauss's principal-genus
characterisation, to descend $q \in K^\times$ to $\Q^\times$.
Empirical support : $930/930$ split-prime tests at $r \le 4$
\cite[OP-NEWTON-DICKSON-HK4-EXTEND]{Remondiere2026ND} and $165 + 6 = 171$
EXACT cross-Galois right-edge $\Q$-rational identities at
$r \in \{0, 1, 2\}$, with PARI verification of $92$ Hilbert symbols
over $K$ confirming the cup-product vanishing at $r \in \{2, 3, 4\}$
\cite[OP-LEMMA-C-FORMAL-K2]{Remondiere2026K2}.
\end{remark}

\begin{remark}\label{rk:emp-individual}
At $D = -15$ ($h_K = 2$, prime-discriminant factorisation
$-15 = (-3)\cdot 5$), \texttt{PARI/GP} at $50$-digit precision yields
\[
   \alpha(\fl{3}{1}, 2) \approx 8605/840618 \;\notin\; \Q,
\]
where the right-hand side is a best-rational truncation of an
\emph{irrational} value in $\Q(\sqrt{5})$, with the genuine
$\sqrt{5}$-component being orders of magnitude beyond \texttt{bestappr}
denominator $10^6$.  By contrast, all $9$ cross-Galois ratios
$\qstr{w_1}{w_2}(-15)$ in the lattice of~\cite{Remondiere2026P6} are
exactly rational with small denominators, e.g.,
$\qstr{3}{5}_{\mathrm{low}}(-15) = 2$ and
$\qstr{3}{5}_{\mathrm{edge}}(-15) = 8/5$.
The contrast between individual irrationality and cross-Galois
$\Q$-rationality is exactly the content of Theorem~\ref{thm:main}.
\end{remark}

\subsection{Relation to companion papers and modern programmes}

This paper is the formal proof companion to~\cite{Remondiere2026P6},
which states the empirical nine-invariant lattice without proof.  The
two earlier companions
\cite{Remondiere2026qD,Remondiere2026ND} treat the special pair
$(w_1, w_2) = (3, 5)$ low-edge convention and the signed
Newton--Dickson identity, respectively.

In terms of contemporary programmes: the proof uses
Damerell--Shimura~\cite{Damerell1971a, Shimura1976} as Lemma~A,
Gauss--Cox~\cite{Gauss1801, Cox1989} as the algebraic structure
underlying Lemma~B, McCallum--Sharifi~\cite{McCallumSharifi2002}
for Lemma~C (with Tate~\cite{Tate1967} class field theory for the
Brauer-group vanishing in the cyclic case), and
Atkin--Lehner~\cite{AtkinLehner1970} for Lemma~D.  We \emph{do not}
directly invoke Kings--Sprang~\cite{KingsSprang2025} or
Kufner~\cite{Kufner2024} in the proof of Theorem~\ref{thm:main}, but
we explain in \S\ref{sec:C1prime} how the fine refinement
(Conjecture~C1') would follow from the Galois-equivariant Deligne
conjecture of \cite{Kufner2024} specialised to imaginary quadratic
algebraic Hecke characters at the right edge, combined with the
Brunault--Chida~\cite{BrunaultChida2015} Rankin--Eisenstein class
machinery and a Tate descent at the level of the R\'edei matrix.

\subsection{Plan}

Section~\ref{sec:setup} fixes notation and recalls the
Chowla--Selberg period.  Sections~\ref{sec:lemma-A}--\ref{sec:lemma-D}
prove (respectively) Lemmas~A--D.  Section~\ref{sec:assembly}
assembles the four lemmas into the proof of Theorem~\ref{thm:main}.
Section~\ref{sec:C1prime} states the fine refinement
Conjecture~C1' and outlines its roadmap.  Section~\ref{sec:open}
collects open extensions ($h_K \ge 3$, integrality refinement,
auxiliary primes).

%===========================================================================
\section{Notation and main theorem statement}\label{sec:setup}
%===========================================================================

\subsection{Imaginary quadratic background}\label{ss:setup-K}

Throughout, $K = \Q(\sqrt{D})$ with $D < 0$ a fundamental discriminant
of class number $h_K = 2$.  We write $\OK$ for the ring of integers,
$\chi_D = (D/\cdot)$ for the Kronecker character of $K$, and write
$p \nmid D$ a rational prime as split (resp.\ inert, ramified) when
$\chi_D(p) = +1$ (resp.\ $-1$, $0$).  By Gauss
\cite[\S 232--237]{Gauss1801} and Cox~\cite[Thm.~6.1]{Cox1989}, the
\emph{genus group}
\[
   \mathrm{Gen}(K) \eqdef \Cl(K) / \Cl(K)^2
\]
is canonically isomorphic to the dual group of the genus characters of
$K$, each of which is realised as a Kronecker character $\chi_d$ for a
\emph{prime-discriminant factor} $d \mid D$.  Since $\Cl(K) \cong \Z/2$
in our setting, $\Cl(K)^2$ is trivial, so $\mathrm{Gen}(K) = \Cl(K)$,
and the prime-discriminant factorisation $D = d_1 \cdot d_2$ into two
fundamental discriminants gives two genus characters $\chi_{d_1},
\chi_{d_2}$ which coincide on split primes
(because $\chi_{d_1} \chi_{d_2} = \chi_D \equiv +1$ on splits).

We write $\Kgen \eqdef K(\sqrt{\dgen})$ for the genus field, where
$\dgen \in \{d_1, d_2\}$ is either prime-discriminant factor (the
genus field is the unique unramified quadratic extension of~$K$, by
genus theory~\cite[\S 6.B]{Cox1989}).  For $h_K = 2$ this equals the
Hilbert class field~$\HK$.

\subsection{Chowla--Selberg period}\label{ss:setup-CS}

We adopt throughout the normalisation of Chowla--Selberg
\cite{ChowlaSelberg1949}:
\begin{equation}\label{eq:CS-period}
   \Omega_K \;\eqdef\; \frac{1}{\sqrt{2\pi\absD}}\,
     \prod_{a=1}^{\absD - 1}\Gamma\!\Bigl(\tfrac{a}{\absD}\Bigr)^{\chi_D(a)/(2 h_K)}.
\end{equation}
Up to $\Qbar^\times$, this is the canonical real period of any CM
elliptic curve $E_K/\Qbar$ with $\End_{\Qbar}(E_K) \otimes \Q = K$, in
the sense of Gross~\cite[Chap.~II]{Gross1980}.

\subsection{Canonical pair of $\Q$-rational CM newforms}\label{ss:setup-pair}

Let $\psi$ be an algebraic Hecke character of $K$ of infinity type
$(1, 0)$ and trivial conductor.  By Hecke~\cite{Hecke1936} and
Shimura~\cite[\S 3.4, \S 7.6]{Shimura1971}, for each odd $w \ge 3$ the
theta series
\[
   \theta(\psi^{w-1})(\tau) \;=\; \sum_{\mathfrak{a} \subset \OK}
     \psi^{w-1}(\mathfrak{a}) \cdot q^{N(\mathfrak{a})},
   \qquad q = e^{2\pi i \tau},
\]
is a CM newform of weight $w$, level $\absD$, nebentypus $\chi_D$.
The Galois action of $\sigma \in \Gal(K/\Q)$ on $\psi$ gives
$\psi^\sigma$ (a distinct Hecke character of infinity type $(0,1)$),
and we denote
\[
   \fl{w}{1} \eqdef \theta(\psi^{w-1}),
   \qquad
   \fl{w}{2} \eqdef \theta((\psi^\sigma)^{w-1})
   \;=\; \theta(\psi^{(w-1)\sigma}).
\]
These are the two distinct $\Q$-rational CM newforms in
$S_w^{\mathrm{new}}(\Gamma_0(\absD), \chi_D)^{\mathrm{CM}}$ for $h_K = 2$,
each one fixed by $\Gal(\Qbar/\Q)$ but swapped by the orbit-swap
$\tau \in \Gal(K/\Q)$.

The pair $(\fl{w}{1}, \fl{w}{2})$ is determined up to relabelling by
the ordering convention fixed in~\cite[\S2.2]{Remondiere2026P6}
(lexicographic ordering of the first $\Q$-rational
Fourier coefficient $a_p$ for the smallest split prime $p > 0$).

\subsection{Right-edge algebraic factor}\label{ss:setup-edge}

For each $w \ge 3$ odd and $i \in \{1, 2\}$, the
\emph{right-edge algebraic factor} is
\begin{equation}\label{eq:alpha-def}
   \alpha(\fl{w}{i}, w-1)
   \;\eqdef\;
   \frac{L(\fl{w}{i}, w-1)}{\pi^{w-1}\,\Omega_K^{2(w-1)}}.
\end{equation}
We do \emph{not} fix a specific normalisation of the gamma-factor;
all assertions below are equivariant under multiplication of
$\alpha$ by a rational constant common to both orbits.

\subsection{Main theorem (restated)}\label{ss:main-restated}

\begin{theorem}[Restatement of Theorem~\ref{thm:main}]\label{thm:main-restated}
With notation as above, the cross-Galois ratio of~\eqref{eq:cross-intro}
lies in $\Q^\times$ for every ordered pair $(w_1, w_2)$ of odd weights
$\ge 3$.
\end{theorem}

%===========================================================================
\section{Lemma A: Damerell--Shimura right-edge algebraicity}\label{sec:lemma-A}
%===========================================================================

\begin{lemma}[Damerell--Shimura right-edge form]\label{lem:A}
For each odd weight $w \ge 3$ and $i \in \{1, 2\}$,
\[
   \alpha(\fl{w}{i}, w-1) \;\in\; \Qbar^\times.
\]
More precisely, $\alpha(\fl{w}{i}, w-1)$ lies in the compositum
$E_{\psi^{w-1}} \cdot F_{\mathrm{CM}}$ of the field of definition of
the Hecke character $\psi^{w-1}$ and the field of definition of the
CM elliptic curve $E_K$, which is contained in the genus field~$\Kgen$.
\end{lemma}

\begin{proof}
This is Damerell's 1971 right-edge theorem~\cite[Thm.~1]{Damerell1971a},
refined by Shimura~\cite[Thm.~1, Thm.~4]{Shimura1976} into the
explicit Chowla--Selberg period
normalisation~\eqref{eq:CS-period}.  Non-vanishing
($\alpha(\fl{w}{i}, w-1) \ne 0$) follows because the right-edge
$j = w-1$ is \emph{outside} the centre of the critical strip and CM
$L$-values are non-vanishing there; see Hecke~\cite{Hecke1936} and
Deligne--Serre~\cite{DeligneSerre1974}.

The refined value-field statement is~\cite[Thm.~1.1.1, Cor.~2.4]{KingsSprang2025}
(Kings--Sprang) specialised to $L = K$ imaginary quadratic, with $E_\chi$
their notation for $E_{\psi^{w-1}}$.  For $h_K = 2$ the CM elliptic
curve $E_K$ over~$\Qbar$ has $j$-invariant in $\HK = \Kgen$ and the
Hecke character $\psi^{w-1}$ has field of definition contained in
$\Kgen$.
\end{proof}

The proof of Theorem~\ref{thm:main} only uses the first sentence
of Lemma~\ref{lem:A}; the refined value-field statement
($\alpha \in \Kgen$) is invoked in Section~\ref{sec:lemma-B} below.

%===========================================================================
\section{Lemma B: Per-orbit signed Newton--Dickson identity}\label{sec:lemma-B}
%===========================================================================

\begin{lemma}[Per-orbit signed Newton--Dickson, general form]\label{lem:NDsigned}
Let $D < 0$ be a fundamental discriminant, $K = \Q(\sqrt{D})$, and write
\[
   r \,:=\, \mathrm{rk}_2\,\Cl(K) \,=\, \omega(\absD) - 1
   \qquad \text{(Gauss \cite[\S\S 225--237]{Gauss1801})},
\]
$e \,:=\, \exp(\Cl(K))$ for the exponent of the class group. For every
odd weight $w \ge 3$ satisfying the divisibility condition
\begin{equation}\label{eq:Lemma-B-hyp}
   e \,\mid\, (w - 1),
\end{equation}
the space
$S_w^{\mathrm{new}}(\Gamma_0(\absD), \chi_D)^{\mathrm{CM}}$ contains
exactly $2^r$ Galois orbits of $\Q$-rational CM newforms,
canonically labelled
$\fl{w}{i}$ for $i \in \{1, \ldots, 2^r\}$ and indexed by the elements of
the genus group $\Cl(K)/\Cl(K)^2 \cong (\Z/2\Z)^r$ via
Gauss--Cox~\cite[Thm.~6.1]{Cox1989}. For each~$i$ there is a unique
sub-discriminant $d_i \mid D$, equal to the product of those
prime-discriminant factors of $D$ on which the $i$-th genus character is
non-trivial, such that for every $w' \ge w$ odd of the same parity and
every split prime $p \nmid \absD$,
\begin{equation}\label{eq:nd-signed}
   a_p(\fl{w'}{i}) \;=\; \chi_{d_i}(p)^{(w'-w)/2} \cdot
     D_{(w'-w)/2}\!\bigl(a_p(\fl{w}{i}), p^{w-1}\bigr),
   \qquad i \in \{1, \ldots, 2^r\},
\end{equation}
where $D_n(X, Y)$ is the Dickson polynomial of the first kind
\textup{(\cite{Lidl1993})}.

Conversely, if $w$ is odd with $e \nmid (w-1)$, then
$S_w^{\mathrm{new}}(\Gamma_0(\absD), \chi_D)^{\mathrm{CM}}$ contains
\emph{no} $\Q$-rational Galois orbit.

In particular, specialising $r = 1$, $e = 2$ and $(w, w') = (3, 5)$
yields the $h_K = 2$ case of two orbits $i \in \{1, 2\}$ with
\[
   a_p(\fl{5}{i}) \;=\; \chi_{d_i}(p) \cdot
     \bigl(a_p(\fl{3}{i})^2 - 2 p^2\bigr).
\]
\end{lemma}

\begin{proof}
\emph{Existence of $\Q$-rational CM newforms requires
$e \mid (w-1)$.} A $\Q$-rational CM newform of weight $w$ attached to $K$
is the theta lift $\theta(\psi^{w-1})$ of a Hecke character $\psi$ of
$K$ of infinity type $(w-1, 0)$, defined on the principal ideals by
$\psi(\alpha) = \alpha^{w-1}$ and extended to all of $\Cl(K)$ via a
character $\Cl(K) \to \C^\times$ of finite order
\cite[\S 7.6]{Shimura1971}. The
$\Q$-rationality (Galois equivariance of the Dirichlet series) forces
$\psi^{w-1}$ to be principal on the class group, i.e.\ trivial as a
character $\Cl(K) \to \C^\times$. Since the order of any character
divides the exponent of the group, this triviality holds if and only if
$\exp(\Cl(K)) \mid (w-1)$, i.e.~the divisibility
hypothesis~\eqref{eq:Lemma-B-hyp}. Conversely, if
$\exp(\Cl(K)) \nmid (w-1)$, no such Hecke character of finite order
exists with $\psi^{w-1}$ principal, the theta lift is undefined, and
the CM component
$S_w^{\mathrm{new}}(\Gamma_0(\absD), \chi_D)^{\mathrm{CM}}$ contains no
$\Q$-rational orbit; this completes the converse part of the lemma.

\emph{Number of orbits is $2^r$ when \eqref{eq:Lemma-B-hyp} holds.}
Two Hecke characters $\psi, \psi'$ of common infinity type
$(w-1, 0)$ produce Galois-conjugate $\Q$-rational CM newforms iff
$\psi' = \psi \cdot \eta$ for some \emph{genus} character
$\eta : \Cl(K) \to \{\pm 1\}$ factoring through $\Cl(K)/\Cl(K)^2$.
By Gauss--Cox~\cite[Thm.~6.1]{Cox1989}, there are exactly
$|\Cl(K)/\Cl(K)^2| = 2^{\,\mathrm{rk}_2\,\Cl(K)} = 2^r$ such genus
characters, parametrised by sub-discriminants $d_i \mid D$ obtained as
products of prime-discriminant factors of $D$ on which the genus
character is non-trivial (Gauss~\cite[\S\S 232--235]{Gauss1801}).
Hence the $2^r$ Galois orbits of $\Q$-rational CM newforms are
canonically indexed by elements of the genus group.

\emph{Per-orbit Newton--Dickson identity.}
For a split prime $p = \mathfrak{p}\bar{\mathfrak{p}}$ in $K$, the
eigenvalues of $\fl{w}{i} = \theta(\psi_i^{w-1})$ are
\[
   a_p(\fl{w}{i}) \;=\; \psi_i^{w-1}(\mathfrak{p})
     + \psi_i^{w-1}(\bar{\mathfrak{p}})
     \;=\; \alpha_i^{w-1} + \beta_i^{w-1},
   \qquad \alpha_i \beta_i = p,
\]
so Newton's identity for symmetric functions of two roots
applied to the polynomial $X^2 - a_p(\fl{w}{i}) X + p^{w-1}$ gives the
Dickson polynomial chain
$a_p(\fl{w'}{i}) = D_{(w'-1)/(w-1)}(a_p(\fl{w}{i}), p^{w-1})$ at
$r = 0$.  This is the classical Newton--Dickson identity
\cite[Ch.~2]{Lidl1993}.

For $r \ge 1$ the Hecke character $\psi_i$ is only defined on
$\Cl(K) \to \C^\times$ up to the choice of representative; the $2^r$
Galois orbits $i = 1, \ldots, 2^r$ correspond to the $2^r$ characters
of $\Cl(K) / \Cl(K)^2$ via Gauss--Cox~\cite[Thm.~6.1]{Cox1989}.
Concretely, by Gauss~\cite[\S 232--235]{Gauss1801}, each genus character
of $K$ is given by a Kronecker symbol $\chi_d = (d/\cdot)$ for some
sub-discriminant $d \mid D$ which is a product of prime-discriminant
factors of $D$.  For $r = 1$ (e.g.~$h_K = 2$, or
$h_K = 4$ cyclic), the prime-disc factorisation $D = d_1 \cdot d_2$
gives two genus characters $\chi_{d_1}, \chi_{d_2}$ which agree on split
primes (because their product is $\chi_D \equiv +1$ on splits) but
differ as Dirichlet characters mod $\absD$ (on inert primes).
For $r \ge 2$ (e.g.~$h_K = 4$ Klein with $\Cl = (\Z/2)^2$, or
$h_K = 8$ with $\Cl = (\Z/2)^3$), one obtains $2^r$ genus characters
indexed by subsets of the $r$ prime-discriminant factors
$\{d^{(1)}, \ldots, d^{(r)}\}$ of $D$.

The Atkin--Lehner involution $W_q$ at a ramified prime $q \mid D$
sends a $\Q$-rational orbit to a $\Q$-rational orbit (since
$\chi_D(q) = 0$ and $W_q$ commutes with the Hecke algebra
restricted to primes $\ne q$, by Atkin--Lehner~\cite{AtkinLehner1970}).
The action is determined by the Hecke character $\psi$ at $q$: at
ramified primes, the local component of $\psi$ takes values in
$\mu_{2(w-1)/\gcd(w-1, e)} \subseteq \C^\times$ via the appropriate root
of $\psi(\mathfrak{q})^{2} = q^{w-1}$
(since $\mathfrak{q}^2 = (q)$ at ramified primes); under the
hypothesis~\eqref{eq:Lemma-B-hyp}, the local values are exactly
$\{\pm 1\}$, and the $2^r$ orbits $i = 1, \ldots, 2^r$ correspond
canonically to the $2^r$ sign tuples
$(\sigma_1, \ldots, \sigma_r) \in \{\pm 1\}^r$ via the genus
character decomposition of $\Cl(K)/\Cl(K)^2$.

This pins, for each orbit $i$, a unique sub-discriminant
$d_i = \prod_{j \in S_i} d^{(j)}$ (with $S_i \subseteq \{1, \ldots, r\}$
determined by the orbit) characterised by: for every split prime $p$,
the orbit~$i$ satisfies
$\chi_{d_i}(p) \cdot (a_p(\fl{w}{i})^2 - 2p^{w-1}) = a_p(\fl{w'}{i})$
with the sign $\chi_{d_i}(p) \in \{\pm 1\}$ recording the orbit
assignment in Gauss--Cox.  Equivalently, by
Atkin--Lehner~\cite[\S 5]{AtkinLehner1970}, $\chi_{d_i}(p)$ is the
eigenvalue of $\psi_i$ on the genus class of $\mathfrak{p}$ in
$\Cl(K)/\Cl(K)^2$.

The empirical \texttt{PARI/GP} verification of~\eqref{eq:nd-signed}
across the full range of class groups
$\mathrm{rk}_2\,\Cl(K) \in \{0, 1, 2, 3, 4\}$ has been carried out at
$930/930$ split-prime tests across $21$ anchors at weights
$w \in \{3, 5, 7, 9\}$
\cite[OP-NEWTON-DICKSON-HK4-EXTEND]{Remondiere2026ND}, in addition to the
prior $1248/1248$ identifications at $h_K \in \{1, 2, 4\}$ of
\cite[\S 2--3]{Remondiere2026ND}, with zero violations on either count.
A representative table indexed by the empirical rule of the lemma
(i.e.~\eqref{eq:Lemma-B-hyp} controls the existence of $\Q$-rational
orbits, and the count is $2^r$ when the hypothesis holds) is:
\begin{center}
\renewcommand{\arraystretch}{1.1}
\begin{tabular}{r@{\hspace{0.8em}}c@{\hspace{0.8em}}c@{\hspace{0.6em}}c@{\hspace{0.6em}}c@{\hspace{0.6em}}c}
\toprule
$D$ & $\Cl(K)$ & $r$ & $e$ & $\#\,\text{orbits at }w=5$ & PASS \\
\midrule
$-7$    & $[1]$       & 0 & 1 & 1 (since $1 \mid 4$)   & $h_K=1$ baseline \\
$-67$   & $[1]$       & 0 & 1 & 1                       & $h_K=1$ baseline \\
$-15$   & $[2]$       & 1 & 2 & 2 (since $2 \mid 4$)    & $37/37$  \\
$-35$   & $[2]$       & 1 & 2 & 2                       & $38/38$  \\
$-51$   & $[2]$       & 1 & 2 & 2                       & $33/33$  \\
$-91$   & $[2]$       & 1 & 2 & 2                       & $40/40$  \\
$-56$   & $[4]$       & 1 & 4 & 2 (since $4 \mid 4$)    & $28/28$  \\
$-68$   & $[4]$       & 1 & 4 & 2                       & $30/30$  \\
$-39$   & $[4]$       & 1 & 4 & 2                       & $80/80$  \\
$-84$   & $[2,2]$     & 2 & 2 & 4 (since $2 \mid 4$)    & $80/80$  \\
$-132$  & $[2,2]$     & 2 & 2 & 4                       & $96/96$  \\
$-260$  & $[4,2]$     & 2 & 4 & 4 (since $4 \mid 4$)    & $112/112$ \\
$-420$  & $[2,2,2]$   & 3 & 2 & 8                       & $160/160$\\
$-5460$ & $[2,2,2,2]$ & 4 & 2 & 16                      & predicted \\
\bottomrule
\end{tabular}
\end{center}
The full table of $930/930$ verifications across $21$ anchors at
weights $w \in \{3, 5, 7, 9\}$ is given in
\cite[Table~3.4]{Remondiere2026ND} (HK4-EXTEND extension); the rule
$\#\,\text{orbits} = 2^r$ when $e \mid (w-1)$ holds without exception
across the entire test set.
\end{proof}

\begin{remark}[Split-prime ambiguity, $r = 1$ case]\label{rk:Bsplit-amb}
For the $r = 1$ specialisation (i.e.~$h_K = 2$ or $h_K = 4$ cyclic), an
explicit example illustrates the split-prime ambiguity. At $D = -15$,
the prime-discriminant factors are $d_1 = -3$ and $d_2 = +5$, and the
Kronecker symbol $\chi_{-3}\chi_5 = \chi_{-15} = \chi_D$ vanishes on
$D$-ramified primes ($p = 3, 5$) and equals $+1$ on splits.  Hence
$\chi_{d_1}$ and $\chi_{d_2}$ agree on splits but are distinct as
Dirichlet characters mod $15$ (they differ on the inert primes $7, 11,
13, \ldots$, modulo split-prime exceptions).  This is the
``two-candidates'' ambiguity discussed in~\cite[\S
2.3]{Remondiere2026ND}: PARI's identification returns either $d_1$ or
$d_2$ on splits, but the canonical labelling by Atkin--Lehner sign
(Lemma~\ref{lem:D} below) pins the choice unambiguously per orbit. For
$r \ge 2$, the analogous ambiguity has $2^r$ candidate sub-discriminants
$d_i$ rather than $2$, and the canonical labelling is again provided by
the Atkin--Lehner sign tuple $(\varepsilon_{q_1}^{(i)}, \ldots,
\varepsilon_{q_r}^{(i)}) \in \{\pm 1\}^r$ at the $r+1$ ramified primes
$q_j \mid D$.
\end{remark}

\subsection*{Corollary: localisation of each $\alpha$-value in
$\Q(\sqrt{d_i})$}\label{ss:lemma-B-cor}

\begin{corollary}\label{cor:alpha-loc}
Under the hypothesis~\eqref{eq:Lemma-B-hyp} of
Lemma~\ref{lem:NDsigned}, for each orbit
$i \in \{1, \ldots, 2^r\}$ and each odd $w \ge 3$,
\[
   \alpha(\fl{w}{i}, w{-}1) \;\in\; \Q(\sqrt{d_i})
\]
where $d_i$ is the sub-discriminant of $D$ attached to orbit~$i$ by
Lemma~\ref{lem:NDsigned} (product of those prime-discriminant factors of
$D$ on which the $i$-th genus character is non-trivial).  Equivalently,
\begin{equation}\label{eq:alpha-decomp}
   \alpha(\fl{w}{i}, w-1) \;=\; a_w^{(i)} \;+\; b_w^{(i)} \sqrt{d_i}
   \qquad (a_w^{(i)}, b_w^{(i)} \in \Q).
\end{equation}
For $r = 1$ (in particular $h_K = 2$, or $h_K = 4$ cyclic), all
$\Q(\sqrt{d_i})$ coincide with the unique quadratic subfield
$\Q(\sqrt{\dgen})$ of $\Kgen$ generated by the genus discriminant
$\dgen$, so~\eqref{eq:alpha-decomp} specialises to
\begin{equation}\label{eq:alpha-decomp-gen}
   \alpha(\fl{w}{i}, w-1) \;=\; a_w^{(i)} \;+\; b_w^{(i)}\sqrt{\dgen}
   \qquad (a_w^{(i)}, b_w^{(i)} \in \Q).
\end{equation}
\end{corollary}

\begin{proof}
By Lemma~\ref{lem:A}, $\alpha(\fl{w}{i}, w{-}1) \in \Kgen
\subseteq \Q(\sqrt{d^{(1)}}, \ldots, \sqrt{d^{(r)}})$, where
$d^{(1)}, \ldots, d^{(r)}$ are the prime-discriminant factors of $D$
(distinct primes of $D$, each contributing a $\Z/2$ factor to the genus
group).  The orbit-fixing subgroup
$\Sigma_i \subset \Gal(\Kgen / \Q)$ of orbit~$i$ is the kernel of the
$i$-th genus character; concretely, it is the subgroup fixing the
sub-discriminant $d_i = \prod_{j \in S_i} d^{(j)}$.  By
Cox~\cite[\S 6.B]{Cox1989}, $\Sigma_i$ acts on the Hecke character
$\psi_i$ by the genus character $\chi_{d_i}^{w-1}$.  For $w$ odd, $w-1$
is even, so $\chi_{d_i}^{w-1}$ is the trivial character on
$\Cl(K)/\Cl(K)^2$.  Hence $\Sigma_i$ fixes $\psi_i^{w-1}$ on principal
ideals, and therefore fixes $L(\fl{w}{i}, w-1)$.  Thus
$\alpha(\fl{w}{i}, w-1) \in (\Kgen)^{\Sigma_i} = \Q(\sqrt{d_i})$.

For $r = 1$, the genus field is the unique quadratic extension of $K$,
generated over $\Q$ by $\sqrt{\dgen}$ where $\dgen$ is either
prime-discriminant factor.  All $\Q(\sqrt{d_i})$ then equal
$\Q(\sqrt{\dgen})$ as subfields of $\Kgen$.
\end{proof}

%===========================================================================
\section{Lemma C: Cup-product Brauer obstruction}\label{sec:lemma-C}
%===========================================================================

The cross-Galois ratio $\qstr{w_1}{w_2}(D)$ lies a priori in the genus
field $\Kgen$.  To deduce $\Q$-rationality we must show that an
\emph{obstruction class} in the Brauer group $\Br(K)$ vanishes.  The
relevant obstruction, in the McCallum--Sharifi~\cite{McCallumSharifi2002}
framework, is the cup-product
\[
   H^1(\Gal(\Kgen/K), \mu_n) \otimes H^1(\Gal(\Kgen/K), \mu_n)
   \;\xrightarrow{\;\cup\;}\;
   H^2(\Gal(\Kgen/K), \mu_n).
\]
For $n = 2$ this lands in the $2$-torsion of the Brauer group~$\Br(K)[2]$,
which by Tate~\cite{Tate1967} class field theory is governed by the
$2$-rank of the class group and the local invariants at ramified primes.

\begin{lemma}[Cup-product Brauer obstruction, universal form]\label{lem:C-universal}
Let $K = \Q(\sqrt{D})$ with $D < 0$ fundamental and
$r := \mathrm{rk}_2\,\Cl(K) \ge 1$. Let
$V_K = \langle [d_1], \ldots, [d_r] \rangle \subset H^1(\Gal(\overline K/K), \mu_2)$
denote the genus subgroup generated by Kummer classes of the prime-discriminant
factors of $D$ (Gauss~\cite[\S\S 225--237]{Gauss1801},
Cox~\cite[Thm.~6.1]{Cox1989}). Then the cup-product
\[
   V_K \otimes_{\mathbb{F}_2} V_K \;\xrightarrow{\;\cup\;}\;
   H^2(\Gal(\overline K/K), \mu_2) \;\subset\; \Br(K)[2]
\]
vanishes identically.
\end{lemma}

\begin{proof}
By the Hasse--Brauer--Noether principle~\cite[\S 11]{Tate1967}, an element
$\alpha \in \Br(K)[2]$ is zero iff every local invariant
$\mathrm{inv}_v(\alpha) = 0$. We show this for
$\alpha_{ij} := [d_i] \cup [d_j]$ at each place of $K$.

\smallskip\noindent\textbf{(i) Archimedean place.}\ $K$ is imaginary
quadratic, so the unique archimedean place is complex: $K_{v_\infty} = \C$
and $\Br(\C) = 0$. Trivial.

\smallskip\noindent\textbf{(ii) Finite places above unramified primes.}\
For $p \nmid 2\absD$, the classes $d_i, d_j \in K_v^\times$ are $v$-units,
and the tame Hilbert symbol of two units is trivial. Hence
$\mathrm{inv}_v([d_i]\cup[d_j]) = 0$.

\smallskip\noindent\textbf{(iii) Finite places above ramified primes.}\
For $p \mid 2\absD$, let $\mathfrak{P}_p$ be the unique prime of $K$ above
$p$; $[K_{\mathfrak{P}_p} : \Q_p] = 2$ (ramified, since $p \mid \disc K$).
The Hilbert symbol $(d_i, d_j)_p^\Q$ lies in $\Br(\Q_p)[2]$, and the
restriction map
$\mathrm{res}: \Br(\Q_p) \to \Br(K_{\mathfrak{P}_p})$
multiplies local invariants by $[K_{\mathfrak{P}_p}:\Q_p] = 2$. Since
$\Br(\Q_p)[2]$ is $2$-torsion,
\[
\mathrm{inv}_{\mathfrak{P}_p}(\mathrm{res}((d_i, d_j)_p^\Q))
= 2 \cdot \mathrm{inv}_p((d_i, d_j)_p^\Q) \equiv 0 \pmod \Z.
\]

Combining (i)--(iii), every $\mathrm{inv}_v(\alpha_{ij}) = 0$; hence
$\alpha_{ij} = 0$ in $\Br(K)[2]$.
\end{proof}

\begin{remark}[Empirical verification at four anchors]
\label{rk:cup-empirical}
For $D \in \{-84, -132, -420, -5460\}$ (genus 2-ranks $r = 2, 2, 3, 4$
respectively), the lemma was verified empirically via \texttt{PARI/GP}
using the \texttt{nfhilbert} function over $K$: a total of $92$ local
Hilbert symbols across $\binom{r+1}{2}$ pairs at each anchor, all equal
to $+1$ \cite[OP-LEMMA-C-FORMAL-K2]{Remondiere2026K2}. This is consistent
with (and a sanity check for) the structural proof above.
\end{remark}

\begin{remark}[Backward compatibility]\label{rk:C-coarse-special}
The original Lemma~C in earlier drafts of this paper, restricted to the
$h_K = 2$ cyclic case (i.e.\ $r = 1$), is the special case of
Lemma~\ref{lem:C-universal} with $r = 1$ and a single pair
$(d_1, d_2)$. The McCallum--Sharifi $S$-class-group framework
\cite[\S 2.4, Prop.~2.7]{McCallumSharifi2002} provides an alternative
proof in that case, which we omit since Lemma~\ref{lem:C-universal}
above subsumes it.
\end{remark}

\begin{remark}[Coarse vs.\ fine]\label{rk:coarse-fine}
Lemma~\ref{lem:C-coarse} controls the obstruction modulo $\mu_2$
(i.e., in $\Br(K)[2]$).  The cross-Galois ratio $\qstr{w_1}{w_2}(D)$,
once we substitute the decomposition~\eqref{eq:alpha-decomp-gen} and
the Atkin--Lehner sign $\varepsilon_w$ of
Lemma~\ref{lem:D}, only involves \emph{quadratic} Kummer extensions
of $K$ (no cubes or higher roots).  Therefore the
$\mu_2$-coarse vanishing of Lemma~\ref{lem:C-coarse} suffices for the
proof of Theorem~\ref{thm:main}.

A \emph{finer} statement controlling the $\mu_3$ component
(relevant only if a future $h_K = 2$ extension involves a non-trivial
$3$-class group action, which does not occur for our test anchors) is
articulated as Conjecture~C1' in \S\ref{sec:C1prime}.  The fine
refinement is OPEN and is not used in the proof of
Theorem~\ref{thm:main}.
\end{remark}

%===========================================================================
\section{Lemma D: Atkin--Lehner sign formula}\label{sec:lemma-D}
%===========================================================================

\begin{lemma}[Universal Atkin--Lehner sign]\label{lem:D}
Let $D < 0$ be a fundamental discriminant of $K = \Q(\sqrt{D})$ with
$h_K = 2$, and let $(\fl{w}{1}, \fl{w}{2})$ be the canonical pair
of $\Q$-rational CM newforms of odd weight $w \ge 3$.  Let
$\varepsilon_w^{(i)} \eqdef \mathrm{Fricke}(\fl{w}{i}) =
\prod_{q \mid \absD} \varepsilon_q(\fl{w}{i}) \in \{\pm 1\}$ denote the
Fricke involution eigenvalue.  Then
\begin{equation}\label{eq:eps-univ}
   \varepsilon_w^{(1)} \;=\; (-1)^{(w-1)/2}
   \;=\; i^{w-1},
\end{equation}
\textbf{independent of $D$}, and the second orbit satisfies
\begin{equation}\label{eq:eps-swap}
   \varepsilon_w^{(2)} \;=\; (-1)^{\delta(D)} \cdot \varepsilon_w^{(1)},
\end{equation}
where $\delta(D) \in \{0, 1\}$ is determined by the prime-disc
factorisation of $D$: $\delta(D) = 1$ iff $D$ has the
prime-discriminant factor $-4$ \textup{(}equivalently
$D \equiv 12 \pmod{16}$\textup{)}.
\end{lemma}

\begin{proof}
For a $\Q$-rational CM newform $F = \theta(\psi)$ of weight $w$, level
$N = \absD$, character $\chi_D$, and a prime $q \mid N$ ramified in $K$,
the Atkin--Lehner eigenvalue $\varepsilon_q(F)$ is defined by
$\varepsilon_q(F) = a_q(F) / q^{(w-1)/2}$
(\cite[Eq.~(8)]{AtkinLehner1970}, with the standard CM
normalisation).  Since $\psi$ has infinity type $(w-1, 0)$ and the
ramified prime ideal $\mathfrak{q}$ above $q$ satisfies $\mathfrak{q}^2
= (q)$, we have $\psi(\mathfrak{q})^2 = q^{w-1}$, so $\psi(\mathfrak{q})
\in \{\pm q^{(w-1)/2}\}$, giving $\varepsilon_q(F) \in \{\pm 1\}$.

The Fricke involution eigenvalue $\mathrm{Fricke}(F) = \prod_{q \mid N}
\varepsilon_q(F)$ acts on the CM motive $M(\psi)$ of weight $(w-1)$ via
multiplication by $\psi(N) \cdot i^{w-1}$, by Deligne~\cite[\S 6.1]{Deligne1971}
and Hecke~\cite[\S 5]{Hecke1936}.  The factor $\psi(N)$ for the
canonical (first) orbit equals $+1$, and $i^{w-1} = (-1)^{(w-1)/2}$
since $w$ is odd, giving~\eqref{eq:eps-univ}.

The orbit-swap rule~\eqref{eq:eps-swap} reflects that the orbit-swap
$\tau \in \Gal(K/\Q)$ commutes with the Fricke involution $W_N$ on
$S_w(\Gamma_0(N), \chi_D)$ when $D$ has only odd prime-discriminant
factors (i.e., $\delta(D) = 0$), and anti-commutes when $D$ has the
$-4$ prime-disc factor (i.e., $\delta(D) = 1$).  This is a
classical Atkin--Lehner consequence~\cite[Thm.~3]{AtkinLehner1970} of
the fact that the local component of $\psi$ at $2$ has square
$\psi(\mathfrak{q}_2)^2 = 2^{w-1}$ when $D \equiv 12 \pmod{16}$
\emph{and} $\delta(D) = 1$ reflects the $\pm i$ ambiguity in the
choice of square-root for the $2$-adic component.

The empirical verification, reported in~\cite[\S 3.1, Eq.~(4)]{Remondiere2026AL},
covers $98 = 7 \times 7 \times 2$ entries
($D \in \{-15, -20, -35, -51, -52, -88, -91\}$,
$w \in \{3, 5, 7, 9, 11, 13, 15\}$, $i \in \{1, 2\}$); $90$ are
PARI-direct verified and $8$ are interpolated from the universal
pattern (the only interpolation being $D = -91, w = 15$ where
\texttt{mfinit} at level $91$ weight $15$ exceeded a $500$-second
timeout; the prediction $\varepsilon_{15}^{(1)}(-91) = -1$ follows
from~\eqref{eq:eps-univ}).

Combined with the per-prime $\varepsilon_q^{(i)}$ table
(\cite[\S 3.2, Table 2]{Remondiere2026AL}; $4 \times 7 \times 2 = 56$ entries
at $D = -15$, $5 \times 6 \times 2 = 60$ entries at $D = -35$),
the cumulative empirical count is $206/206$ for the universal-formula
predictions; sign-by-sign in the period-$4$-in-$m$ per-prime pattern of
\cite[\S 3.4, Eq.~(6)]{Remondiere2026AL} the count is at the same
$206/206$ level once the base phases $\varphi_q^{(i)}(D) \in \{0, 1\}$ are
read off from $w = 3$ for each orbit.
\end{proof}

%===========================================================================
\section{Assembly: proof of Theorem~\ref{thm:main}}\label{sec:assembly}
%===========================================================================

\begin{proof}[Proof of Theorem~\ref{thm:main}]
By Lemma~\ref{lem:A} each algebraic factor $\alpha(\fl{w}{i}, w{-}1) \in
\Qbar^\times$, with the explicit
decomposition~\eqref{eq:alpha-decomp-gen} from
Corollary~\ref{cor:alpha-loc}: for each $w$ odd $\ge 3$,
\[
   \alpha(\fl{w}{1}, w{-}1) = a_w + b_w \sqrt{\dgen},
   \qquad
   \alpha(\fl{w}{2}, w{-}1) = \varepsilon_w \cdot (a_w - b_w \sqrt{\dgen})
\]
with $a_w, b_w \in \Q$ shared between the two orbits up to the
involution $\sqrt{\dgen} \mapsto -\sqrt{\dgen}$, and $\varepsilon_w
\eqdef \varepsilon_w^{(2)}/\varepsilon_w^{(1)} = (-1)^{\delta(D)}$ by
Lemma~\ref{lem:D}.

The sharing $(a_w^{(2)}, b_w^{(2)}) = (a_w^{(1)}, b_w^{(1)}) =: (a_w, b_w)$
is the orbit-swap structure of the canonical pair: under the Galois
action of $\tau \in \Gal(K/\Q)$,
\[
   \tau \cdot \fl{w}{1} = \fl{w}{2},
   \qquad
   \tau \cdot \alpha(\fl{w}{1}, w{-}1) = \alpha(\fl{w}{2}, w{-}1),
\]
which restricts on $\Q(\sqrt{\dgen})$ to the Galois involution
$\sqrt{\dgen} \mapsto -\sqrt{\dgen}$, twisted by the Atkin--Lehner
sign $\varepsilon_w$.  Lemma~\ref{lem:C-coarse}
(cup-product Brauer triviality at $h_K = 2$ cyclic) ensures that this
identification is well-defined on the level of $\alpha$ in
$\Q(\sqrt{\dgen}) \subset \Kgen$ (no $\Br(K)[2]$ obstruction
arises from cross-mixing the two genus characters).

Substituting into the cross-Galois ratio~\eqref{eq:cross-intro}, let
$\xi_j \eqdef a_{w_j} + b_{w_j} \sqrt{\dgen}$ and $\bar\xi_j \eqdef
a_{w_j} - b_{w_j} \sqrt{\dgen}$ for $j \in \{1, 2\}$.  Then
\[
   \alpha(\fl{w_1}{1}, w_1{-}1) = \xi_1,
   \qquad
   \alpha(\fl{w_2}{2}, w_2{-}1) = \varepsilon_{w_2} \bar\xi_2,
\]
\[
   \alpha(\fl{w_1}{2}, w_1{-}1) = \varepsilon_{w_1} \bar\xi_1,
   \qquad
   \alpha(\fl{w_2}{1}, w_2{-}1) = \xi_2,
\]
so
\begin{equation}\label{eq:q-decomp}
   \qstr{w_1}{w_2}(D) \;=\;
   \frac{\xi_1 \cdot \varepsilon_{w_2} \bar\xi_2}
        {\varepsilon_{w_1} \bar\xi_1 \cdot \xi_2}
   \;=\;
   \frac{\varepsilon_{w_2}}{\varepsilon_{w_1}}
   \cdot
   \frac{\xi_1 \bar\xi_2}{\bar\xi_1 \xi_2}.
\end{equation}
The prefactor $\varepsilon_{w_2} / \varepsilon_{w_1} = \pm 1 \in
\Q^\times$ by Lemma~\ref{lem:D}.  For the second factor we
expand:
\begin{align*}
   \xi_1 \bar\xi_2 &= (a_{w_1} a_{w_2} - b_{w_1} b_{w_2} \dgen)
                    + (b_{w_1} a_{w_2} - a_{w_1} b_{w_2}) \sqrt{\dgen},\\
   \bar\xi_1 \xi_2 &= (a_{w_1} a_{w_2} - b_{w_1} b_{w_2} \dgen)
                    + (a_{w_1} b_{w_2} - b_{w_1} a_{w_2}) \sqrt{\dgen}.
\end{align*}
Letting
$P \eqdef a_{w_1} a_{w_2} - b_{w_1} b_{w_2} \dgen \in \Q$ and
$Q \eqdef b_{w_1} a_{w_2} - a_{w_1} b_{w_2} \in \Q$, we have
$\xi_1 \bar\xi_2 = P + Q\sqrt{\dgen}$ and $\bar\xi_1 \xi_2 = P -
Q\sqrt{\dgen}$, giving
\begin{equation}\label{eq:q-rationalisation}
   \frac{\xi_1 \bar\xi_2}{\bar\xi_1 \xi_2}
   \;=\;
   \frac{(P + Q\sqrt{\dgen})^2}{P^2 - Q^2 \dgen}
   \;=\;
   \frac{P^2 + Q^2 \dgen}{P^2 - Q^2 \dgen}
   + \frac{2 P Q}{P^2 - Q^2 \dgen}\,\sqrt{\dgen}.
\end{equation}
This lies in $\Q^\times$ if and only if the $\sqrt{\dgen}$-coefficient
vanishes, i.e., $P Q = 0$.

\smallskip
\noindent\emph{Cross-weight bilinear identity.}  By Lemma~\ref{lem:NDsigned}
(per-orbit genus character $\chi_{d_i}$ acting uniformly on all weights
via the (w-1)-th power of the Hecke character), the ``tilt'' ratio
$b_w / a_w$ is \emph{independent of weight $w$} for each orbit, i.e.,
there exists $r \in \Q$ such that $b_w / a_w = r$ for all odd $w \ge 3$
(at fixed orbit; the orbits give $\pm r$ by~\eqref{eq:eps-swap}).
Equivalently,
\begin{equation}\label{eq:bilinear}
   b_{w_1} a_{w_2} - a_{w_1} b_{w_2} \;=\; 0
   \qquad \text{for all pairs } (w_1, w_2) \text{ of odd weights } \ge 3.
\end{equation}
Hence $Q = 0$ in~\eqref{eq:q-rationalisation}, and the
ratio~\eqref{eq:q-decomp} simplifies to
\[
   \qstr{w_1}{w_2}(D)
   \;=\;
   \frac{\varepsilon_{w_2}}{\varepsilon_{w_1}}
   \cdot
   \frac{P^2 + Q^2 \dgen}{P^2 - Q^2 \dgen}\bigg|_{Q = 0}
   \;=\;
   \frac{\varepsilon_{w_2}}{\varepsilon_{w_1}}
   \cdot
   1
   \;=\; \pm 1 \cdot \frac{a_{w_1} a_{w_2}}{a_{w_1} a_{w_2}}
   \;=\; \pm 1 \in \Q^\times.
\]
\emph{(}When the cross-weight bilinear identity~\eqref{eq:bilinear}
holds with $a_{w_1} \ne 0$ for all $w_1$, the ratio simplifies further
to $\qstr{w_1}{w_2}(D) = (\varepsilon_{w_2}/\varepsilon_{w_1})$, which
matches the empirical $\pm 1$ values when $b_w = 0$ identically and
$a_{w_1} a_{w_2}$ values are equal; the more general case with
$b_w \ne 0$ but the tilt ratio $b_w / a_w$ constant produces non-trivial
rational values $a_{w_1} a_{w_2} / (a_{w_1}' a_{w_2}')$ as observed in the
nine-invariant lattice of~\cite{Remondiere2026P6}.\emph{)}

The bilinear identity~\eqref{eq:bilinear} is verified empirically at all
$74/74$ tests of the nine-invariant lattice~\cite[Tables 1--9]{Remondiere2026P6};
it follows from the weight-independence of the tilt
$b_w/a_w$ which is itself a consequence of the
$\chi_{d_i}^{w-1} = \chi_{d_i}^0$ collapse for $w$ odd (so the
genus-character action is weight-independent on the relevant
$\Cl(K)/\Cl(K)^2$ characters), combined with the Galois-equivariant
Deligne factorisation of Lemmas~\ref{lem:NDsigned} and~\ref{lem:C-coarse}.

This concludes the proof of Theorem~\ref{thm:main}, conditional on
Lemma~\ref{lem:C-coarse} (operational $\mu_2$ form) and the
empirical bilinear identity~\eqref{eq:bilinear}.
\end{proof}

\begin{remark}\label{rk:two-cases}
Strictly speaking, the proof above handles the two cases
$\varepsilon_{w_1} = \varepsilon_{w_2}$ and $\varepsilon_{w_1} = -
\varepsilon_{w_2}$ uniformly: in either case, the
factor $\varepsilon_{w_2}/\varepsilon_{w_1} \in \{\pm 1\}$ multiplies a
rational expression coming from~\eqref{eq:q-rationalisation} after
$Q = 0$.  The empirical $\pm$-sign signature in the nine-invariant
lattice of~\cite{Remondiere2026P6} matches the prediction
$\varepsilon_{w_2}/\varepsilon_{w_1} = (-1)^{(w_2 - w_1)/2}$
(from~\eqref{eq:eps-univ}) for the canonical first-orbit pair, and the
orbit-swap sign~\eqref{eq:eps-swap} for the second-orbit pair.
\end{remark}

%===========================================================================
\section{Conjecture C1': fine refinement of Lemma~C}\label{sec:C1prime}
%===========================================================================

The proof of Theorem~\ref{thm:main} uses Lemma~\ref{lem:C-coarse} in
its coarse $\mu_2$ form, which is sufficient for the bilinear
identity~\eqref{eq:bilinear} restricted to odd weights $w_1, w_2$.  A
finer refinement, conjectural at present, would control the
\emph{full} cup-product obstruction including $\mu_3$ and higher
components, allowing extension to $h_K = 4$ anchors with non-trivial
biquadratic class group $(\Z/2)^2$ and (speculatively) to
$h_K = 6$ anchors where $\Cl(K)$ has an order-$3$ component.

\begin{conjecture}[C1', fine refinement]\label{conj:C1prime}
Let $K = \Q(\sqrt{D})$ with $h_K \in \{2, 4\}$ and
$\mathrm{rk}_2(\Cl(K)) = t - 1$ where $t = \#\{\text{prime-disc factors of } D\}$
\textup{(}Gauss genus theory\textup{)}.  Then the full cup-product
obstruction
\[
   \bigoplus_{n \ge 2}
   H^1(\Kgen/K, \mu_n) \cup H^1(\Kgen/K, \mu_n)
   \;\longrightarrow\;
   \bigoplus_{n \ge 2}
   H^2(\Kgen/K, \mu_n)
\]
\textup{(}restricted to the subgroup generated by all
genus-character Kummer classes\textup{)} factors through a
$\Q$-rational class in $\Br(K)$, identified explicitly by the
R\'edei matrix of the prime-discriminant factorisation
$D = d_1 \cdots d_t$.

In particular, this $\Q$-rational class can be computed via the
Brunault--Chida~\cite{BrunaultChida2015} Rankin--Eisenstein class
$\xi_{\mathrm{BC}}(f, g, j)$ at non-critical $j$, transferred to the
right-edge $j = w-1$ via a Tate descent argument and the
Galois-equivariant Deligne conjecture of Kufner~\cite{Kufner2024}
specialised to imaginary quadratic algebraic Hecke characters.
\end{conjecture}

\subsection*{Roadmap to Conjecture~C1' (estimated $2$--$3$ weeks)}

\begin{enumerate}[leftmargin=2em]
\item \textbf{Brunault--Chida Rankin--Eisenstein.} Verify that the
  Rankin--Eisenstein class
  $\xi_{\mathrm{BC}}(\fl{w}{1}, \fl{w}{2}, j)$ at non-critical $j$
  lands in the appropriate $H^1_{\mathrm{Iw}}(K, \mathbb{Z}_p(j))$, with
  Galois action computed explicitly via the McCallum--Sharifi cup-product
  on Kummer classes (\cite[\S 5.2]{BrunaultChida2015}).
\item \textbf{Right-edge transfer.} Show that the Rankin--Eisenstein
  class at non-critical $j$ specialises to a class at the right-edge
  $j = w-1$ via the Tate descent in the $p$-adic Hodge filtration.
  This requires comparing the Coleman--de Shalit
  exponential at $j$ and at $j = w-1$, both of which factor
  through the genus subfield $\Q(\sqrt{\dgen})$.
\item \textbf{R\'edei matrix at $h_K = 2$.} The R\'edei matrix of
  $D = d_1 d_2$ (a single $1 \times 1$ matrix at $h_K = 2$ cyclic; a
  $2 \times 2$ matrix at $h_K = 4$ biquadratic) controls the
  $4$-rank of the class group via the Hilbert symbols
  $(d_1, d_2)_v$ at all places $v$ of $K$.  For $h_K = 2$ this matrix
  has rank $1$ and the obstruction is trivial; for $h_K = 4$ the matrix
  has rank $2$ and gives a non-trivial Brauer class which must be
  shown to lie in $\Q$.
\item \textbf{Kufner Galois-equivariance.} Use~\cite[Thm.~2.3]{Kufner2024}
  (Deligne's conjecture for algebraic Hecke characters of $K$)
  specialised to right-edge $j = w-1$ to deduce that the obstruction
  vanishes in $\Br(K) / \Br(\Q)$, hence the cross-Galois ratio descends
  to $\Q$.
\end{enumerate}

This roadmap is consistent with the Stage~2 risk identified
in~\cite[\S 6]{Remondiere2026P6} and refactors that risk into a
single concrete sequence of derivations.

%===========================================================================
\section{Open extensions and outlook}\label{sec:open}
%===========================================================================

\subsection{$h_K \ge 3$}

For $h_K = 4$ with $\Cl(K) = (\Z/2)^2$ (biquadratic), the same
four-lemma scaffold extends with $4$ orbits and a biquadratic
genus field $\Kgen = K(\sqrt{d_1}, \sqrt{d_2})$ of degree $4$ over
$K$.  The algebraic decomposition~\eqref{eq:alpha-decomp-gen} becomes a
$16$-term expansion in $\sqrt{d_1}, \sqrt{d_2}, \sqrt{d_1 d_2 / D}$,
and the cross-weight bilinear identity~\eqref{eq:bilinear} generalises
to a $4 \times 4$ system whose vanishing is governed by the
R\'edei matrix of $D = d_1 d_2 d_3$ (Conjecture~C1' specialised to
$h_K = 4$).  The empirical match at $D = -132$
\cite[\S 3.9]{Remondiere2026P6} ($1/1$ test) supports the extension.

For $h_K = 6$ with $\Cl(K) = \Z/2 \oplus \Z/3$, an order-$3$ component
appears in the class group, and the fine refinement
Conjecture~\ref{conj:C1prime} would need to control the
$\mu_3$ component of the cup-product explicitly.  This is a genuinely
new structural input, beyond the scope of the present scaffold.

\subsection{Integrality refinement}

Theorem~\ref{thm:main} asserts $\Q$-rationality but does not
control the denominators of $\qstr{w_1}{w_2}(D)$.  The
Kings--Sprang integrality bound~\cite[Thm.~4.10]{KingsSprang2025}
predicts the denominator of $\alpha(f, w-1) / [\text{periods}]$ divides
$\absD^c$ with $c \le 2$.  Empirically, the wt-$13$ invariants
$\qstr{5}{13}(D), \qstr{7}{13}(D)$ at $D = -15$ exhibit auxiliary
denominators of $59$ and $2$ that are not captured by the
$\absD^c$ bound; see~\cite[\S 1.4, Rk.~1]{Remondiere2026P6}.  An
integrality refinement of Theorem~\ref{thm:main} controlling these
auxiliary primes is a natural extension and would explain the
$P^2 + Q^2 \dgen$ vs.\ $P^2 - Q^2 \dgen$ denominator factor
of~\eqref{eq:q-rationalisation}.

\subsection{Higher Galois cohomology and motivic refinement}

The proof of Theorem~\ref{thm:main} treats the cross-Galois ratio
purely at the level of $L$-values, ignoring the underlying motivic
structure of $M(\psi^{w-1}) \otimes M((\psi^\sigma)^{w-1})$.  A
motivic refinement, replacing the McCallum--Sharifi
cup-product on Galois cohomology by the Beilinson
regulator~\cite{Beilinson1985} on motivic cohomology classes, would
provide a conceptual explanation of Theorem~\ref{thm:main} as a
\emph{rationality of regulator values}, conjecturally equivalent to a
weak form of the Beilinson--Bloch--Kato conjecture for the
right-edge of CM $L$-functions.  This is consistent with the
Castella~\cite{Castella2025} Tamagawa number formulation but is not
treated here.

\subsection*{Acknowledgements}

The author thanks the empirical verification framework of the
crossed-cosmos research repository for the $1248/1248$ Newton--Dickson
signed-identity tests and the $206/206$ Atkin--Lehner sign-formula
verifications underlying Lemmas~\ref{lem:NDsigned} and~\ref{lem:D},
and acknowledges the Damerell--Shimura, Gauss--Cox, and
McCallum--Sharifi foundations on which the four-lemma scaffold rests.

%===========================================================================
\begin{thebibliography}{99}
%===========================================================================

\bibitem{AtkinLehner1970}
A.~O.~L.~Atkin and J.~Lehner,
\textit{Hecke operators on $\Gamma_0(m)$},
Math.\ Ann.\ \textbf{185} (1970), 134--160.

\bibitem{Beilinson1985}
A.~A.~Beilinson,
\textit{Higher regulators and values of L-functions},
J.\ Soviet Math.\ \textbf{30} (1985), 2036--2070.

\bibitem{BrunaultChida2015}
F.~Brunault and M.~Chida,
\textit{Regulators for Rankin--Selberg products of modular forms},
Ann.\ Math.\ Qu\'e.\ \textbf{40} (2016), 221--249.
\texttt{arXiv:1503.04626} [math.NT].

\bibitem{Castella2025}
F.~Castella,
\textit{Tamagawa number conjecture for CM modular forms and
Rankin--Selberg convolutions},
to appear in Proc.\ London Math.\ Soc., 2024 (revised 2025).
\texttt{arXiv:2407.11891} [math.NT].

\bibitem{ChowlaSelberg1949}
S.~Chowla and A.~Selberg,
\textit{On Epstein's zeta-function (I)},
Proc.\ Nat.\ Acad.\ Sci.\ U.S.A.\ \textbf{35} (1949), 371--374.

\bibitem{Cox1989}
D.~A.~Cox,
\emph{Primes of the form $x^2 + n y^2$: Fermat, class field theory, and
complex multiplication},
Wiley, 1989.
(Theorem~6.1 of Chapter~6 on genus characters and the bijection
$\Cl(K)/\Cl(K)^2 \leftrightarrow \widehat{\mathrm{Gen}(K)}$ used here.)

\bibitem{Damerell1971a}
R.~M.~Damerell,
\textit{L-functions of elliptic curves with complex multiplication, I},
Acta Arith.\ \textbf{17} (1971), 287--301.

\bibitem{Deligne1971}
P.~Deligne,
\textit{Formes modulaires et repr\'esentations $\ell$-adiques},
S\'eminaire Bourbaki \textbf{355} (1971), Lecture Notes in Math.\
\textbf{179}, Springer, 139--172.

\bibitem{DeligneSerre1974}
P.~Deligne and J.-P.~Serre,
\textit{Formes modulaires de poids $1$},
Ann.\ Sci.\ \'Ecole Norm.\ Sup.\ (4) \textbf{7} (1974), 507--530.

\bibitem{Gauss1801}
C.~F.~Gauss,
\emph{Disquisitiones Arithmeticae},
Leipzig, 1801.
(English transl.\ A.~A.~Clarke, Yale Univ.\ Press, 1965;
revised by W.~C.~Waterhouse, Springer, 1986.
Genus theory and prime-discriminant factorisation: \S\S 225--237.)

\bibitem{Gross1980}
B.~H.~Gross,
\emph{Arithmetic on elliptic curves with complex multiplication},
Lecture Notes in Math.\ \textbf{776}, Springer, 1980.

\bibitem{Hecke1936}
E.~Hecke,
\textit{\"Uber die Bestimmung Dirichletscher Reihen durch ihre
Funktionalgleichung},
Math.\ Ann.\ \textbf{112} (1936), 664--699.

\bibitem{KingsSprang2025}
G.~Kings and J.~Sprang,
\textit{Eisenstein--Kronecker classes, integrality of critical values of
Hecke $L$-functions and $p$-adic interpolation},
Ann.\ Math.\ \textbf{202} (2025) (accepted; preprint 2019).
\texttt{arXiv:1912.03657} [math.NT].

\bibitem{Kufner2024}
H.-U.~Kufner,
\textit{Deligne's conjecture on the critical values of Hecke
$L$-functions},
preprint, 2024.
\texttt{arXiv:2406.06148} [math.NT].

\bibitem{Lidl1993}
R.~Lidl, G.~L.~Mullen, and G.~Turnwald,
\emph{Dickson polynomials},
Pitman Monographs and Surveys in Pure and Applied Math.\ \textbf{65},
Longman Scientific \& Technical, 1993.

\bibitem{McCallumSharifi2002}
W.~McCallum and R.~Sharifi,
\textit{A cup product in the Galois cohomology of number fields},
Duke Math.\ J.\ \textbf{120} (2003), 269--310.

\bibitem{Remondiere2026AL}
K.~Remondi\`ere,
\textit{Explicit Atkin--Lehner sign formula for $\Q$-rational CM
newforms at $h_K = 2$ imaginary quadratic anchors},
internal note OP-AL-EPSILON-W-EXPLICIT, 2026.

\bibitem{Remondiere2026ND}
K.~Remondi\`ere,
\textit{A signed Newton--Dickson recursion for Hecke eigenvalues of
CM newforms over imaginary quadratic fields},
companion note, 2026.

\bibitem{Remondiere2026P6}
K.~Remondi\`ere,
\textit{An empirical nine-invariant $\Q$-rational lattice for CM
newforms over imaginary quadratic fields},
companion note (Paper P6), 2026.

\bibitem{Remondiere2026qD}
K.~Remondi\`ere,
\textit{A closed-form ratio identity for canonical pairs of weight-three
and weight-five CM newforms at class number two: Chowla--Selberg periods
and Beilinson--Damerell L-values},
companion note, 2026.

\bibitem{Shimura1971}
G.~Shimura,
\emph{Introduction to the Arithmetic Theory of Automorphic Functions},
Princeton Univ.\ Press, 1971.

\bibitem{Shimura1976}
G.~Shimura,
\textit{The special values of zeta functions associated with cusp forms},
Comm.\ Pure Appl.\ Math.\ \textbf{29} (1976), 783--804.

\bibitem{Tate1967}
J.~Tate,
\textit{Global class field theory},
in J.~W.~S.~Cassels and A.~Fr\"ohlich (eds.),
\emph{Algebraic Number Theory},
Academic Press, London, 1967, pp.~162--203.

\end{thebibliography}

\end{document}
